\section{Удиртгал}
\label{sec:introduction}

Багц сонголтын орчин үеийн онолын суурийг \citet{markowitz1952}-ийн дундаж--вариацын загвар тавьсан бөгөөд уг загварт багцын эрсдэлийг хөрөнгийн өгөөжийн ковариацын матрицаар тодорхойлдог. Глобал хамгийн бага вариацтай багц (Global Minimum Variance Portfolio, GMVP) нь хүлээгдэж буй өгөөжийн үнэлгээ шаардахгүйгээр ковариацын матрицад тулгуурлан багцын дисперсийг хамгийн бага байлгах жинг сонгодог. Гэвч жинхэнэ ковариацын матриц үл мэдэгдэх учир түүнийг түүхэн өгөгдлөөс үнэлэх шаардлагатай бөгөөд үнэлгээний алдаа нь багцын жин болон түүврээс гадуурх эрсдэлд нөлөөлж болно.

Түүврийн ковариацын матриц нь $N$ хөрөнгийн хувьд $N(N+1)/2$ ялгаатай параметр агуулдаг. Хөрөнгийн тоо $N$ нь ажиглалтын тоо $T$-тэй харьцуулахад ихсэх үед эдгээр параметрийн үнэлгээ тогтворгүй болох, цаашлаад $N>T$ нөхцөлд түүврийн ковариацын матриц сингуляр болох асуудал үүсдэг \citep{ledoitwolf2004a}. Гэхдээ матриц урвуулагдах боломжтой, $N/T$ харьцаа бага үед ч үнэлгээний алдаа бүрэн арилна гэсэн баталгаа байхгүй. Ялангуяа багцын оновчлол нь ковариацын матрицын үнэлгээний алдааг жингийн сонголтоор дамжуулан өсгөж болзошгүй. \citet{demiguel2009} багц сонголтын олон стратегийг бодит өгөгдөл дээр харьцуулж, параметрийн үнэлгээний алдаа нь оновчлолын практик үр дүнг хязгаарлаж болохыг харуулсан.

Энэ асуудлыг шийдвэрлэх нэг чиглэл нь ковариацын матрицын агшилтын үнэлгээ (shrinkage estimation) юм. Шугаман агшилт нь түүврийн ковариацын матриц $S$-ийг тодорхой бүтцийн таамаглалтай зорилтот матриц $T$-тай $\widehat{\Sigma}=(1-\alpha)S+\alpha T$ хэлбэрээр хослуулдаг бөгөөд $\alpha\in[0,1]$ нь агшилтын эрчим юм. Түүврийн мэдээллийг хадгалах болон үнэлгээний хэлбэлзлийг багасгах хоёр зорилгын тэнцвэр нь зорилтот матрицын бүтэц, агшилтын эрчмээс хамаарна. Тогтмол вариацын зорилтот матриц нь бүх хөрөнгийн дисперсийг ижил, ковариацыг тэг гэж үздэг \citep{ledoitwolf2004a}; диагональ зорилтот матриц нь хөрөнгө тус бүрийн дисперсийг хадгалж, ковариацыг тэглэдэг \citep{ledoit1995}. Тогтмол корреляцын зорилтот матриц нь хөрөнгийн хосуудын корреляцыг нэг нийтлэг утгаар төлөөлүүлдэг \citep{ledoit1995,ledoitwolf2004b}, харин нэг хүчин зүйлийн зорилтот матриц нь өгөөжийн хамтын хөдөлгөөнийг зах зээлийн нийтлэг хүчин зүйлээр загварчилдаг \citep{ledoitwolf2003}. Эдгээр зорилтот матрицын тодорхойлолт, таамаглал болон үнэлгээний томьёог дараагийн онол, арга зүйн хэсэгт дэлгэрэнгүй авч үзнэ.

Шугаман агшилтаас ялгаатай нь шугаман бус агшилт (nonlinear shrinkage, NLS) түүврийн ковариацын матрицын хувийн утгуудыг ялгаатай хэмжээгээр засах боломж олгодог. \citet{ledoitwolf2012} энэ чиглэлийн аргыг хөгжүүлж, \citet{ledoitwolf2020} аналитик хэрэгжилтийг санал болгосон. Түүнчлэн ковариацын матрицын үнэлгээний чанар нь сонгосон алдагдлын функцээс хамаарч болохыг \citet{ledoitwolf2021} судалсан. Иймээс ковариацын матрицыг статистикийн тодорхой шалгуураар сайн үнэлэх нь GMVP-ийн түүврээс гадуурх дисперсийг заавал хамгийн бага болгоно гэсэн үг биш. GMVP-ийн хүрээнд агшилтын аргуудын практик гүйцэтгэлийг судалсан сүүлийн үеийн ажлууд ч зорилтот матриц болон агшилтын эрчмийн сонголт чухал болохыг онцолж байна \citep{shi2025}.

Өмнөх судалгаанууд агшилтын үнэлэгчдийн онолын үндэслэл, өндөр хэмжээст нөхцөл дэх шинж чанар болон багцын оновчлол дахь хэрэглээг өргөн авч үзсэн. Харин хөрөнгийн тоо ажиглалтын тооноос харьцангуй бага, түүврийн ковариацын матриц тооцооллын хувьд ашиглах боломжтой үед агшилтын аргууд GMVP-ийн түүврээс гадуурх эрсдэлийг хэр тогтвортой бууруулах нь тусад нь эмпирик үнэлэх шаардлагатай асуудал хэвээр байна. Ялангуяа нэг багц дээр ажиглагдсан давуу тал нь хөрөнгүүдийн корреляцын бүтэц, багцын бүрэлдэхүүн болон сургалтын цонхны уртаас хамаарах эсэхийг ялган шалгах нь чухал. Иймээс энэ судалгаа агшилтын аргуудын давуу талыг урьдчилан таамаглахын оронд бага $N/T$ нөхцөл дэх тэдгээрийн гүйцэтгэл, тогтвортой байдал болон статистикийн тодорхой бус байдлыг үнэлэхэд чиглэнэ.

Судалгааны зорилго нь түүврийн ковариацын матриц (Sample), тогтмол вариацын (Var), диагональ (Diag), тогтмол корреляцын (Corr), нэг хүчин зүйлийн (SF) зорилтот матрицад суурилсан шугаман агшилтын дөрвөн үнэлэгч болон аналитик шугаман бус агшилтын үнэлэгчийг (NLS) богино борлуулалтыг зөвшөөрөөгүй GMVP-ийн нэг ижил нөхцөлд харьцуулах явдал юм. Үнэлгээнд сарын өгөөж ашиглаж, корреляцын бүтэц ялгаатай, тус бүр 10 хувьцаанаас бүрдсэн өндөр, дунд, бага корреляцтай гурван багцын хувьд 120 сарын сургалтын цонхтой rolling-window аргаар 52 сарын түүврээс гадуурх гүйцэтгэлийг тооцно. Гол үзүүлэлтүүд нь түүврээс гадуурх хэлбэлзэл, Sample аргатай харьцуулсан дисперсийн бууралт (RVR), Sharpe-ийн харьцаа, багцын жингийн эргэлт (turnover) болон төвлөрөл (HHI) юм.

Энэхүү судалгааны эмпирик хувь нэмэр нь хэд хэдэн харилцан нөхөх шинжилгээг нэг хүрээнд хэрэгжүүлсэнд оршино. Нэгдүгээрт, корреляцын ялгаатай бүтэцтэй багцуудын үр дүнг харьцуулж, сургалтын цонхыг 60, 90, 120 сар болгон өөрчлөх замаар $N/T$-ийн мэдрэмжийг үнэлнэ. Хоёрдугаарт, корреляцын ангилал тус бүрээс сонгосон 50, нийт 150 нэмэлт багцад шинжилгээг давтаж, анхны гурван багцын үр дүн бүрэлдэхүүний сонголтоос хэр хамаарч байгааг шалгана. Гуравдугаарт, HAC болон circular moving-block bootstrap аргаар ажиглагдсан ялгааны тодорхой бус байдлыг үнэлж, олон харьцуулалтад Holm-ийн засвар хэрэглэнэ. Дөрөвдүгээрт, тэнцүү жинтэй $1/N$ багцыг нэмэлт жишиг болгон ашиглаж, арилжааны зардлын дараах гүйцэтгэлийг мөн харьцуулна.

Үндсэн үр дүнгээр өндөр корреляцтай багцад Var, дунд корреляцтай багцад Corr үнэлэгч Sample аргатай харьцуулахад түүврээс гадуурх дисперсийг тус тус 5.81\% болон 2.26\%-иар бууруулсан бол бага корреляцтай багцад агшилтын таван үнэлэгчийн RVR бүгд сөрөг байв. Нэмэлт 150 багцын шинжилгээнд үр дүн багцын бүрэлдэхүүнээс хамааран өөрчлөгдсөн. Үндсэн гурван багцад GMVP стратегиуд $1/N$ багцаас бага хэлбэлзэл үзүүлсэн боловч олон харьцуулалтын засварын дараа агшилтын үнэлэгчдийн Sample аргаас, мөн GMVP стратегиудын $1/N$ багцаас статистикийн ач холбогдолтой давуу байдал тогтоогдоогүй. Иймээс бага $N/T$ нөхцөлд агшилтын үр ашгийг зөвхөн үнэлэгчийн онолын шинж чанараар бус, багцын бүтэц болон түүврээс гадуурх гүйцэтгэлд тулгуурлан үнэлэх шаардлагатайг үр дүн харуулж байна.
